\chapter{Introduction}
\label{chap:cz-introduction}

The Hamilton--Perelman Ricci-flow program proves the Poincar\'e
conjecture and Thurston's geometrization conjecture.  The Poincar\'e
conjecture asserts that every closed $3$-manifold with trivial
fundamental group is homeomorphic to $\mathbb S^3$ \cite{Th82}.

For a compact $3$-manifold, evolve an initial metric by
\[
\frac{\partial g_{ij}}{\partial t}=-2R_{ij}
\]
introduced by Hamilton \cite{Ha82}.  The program analyzes finite-time
singularities, performs surgery in
canonical neck and cap regions, and studies the resulting long-time
solution.  A maximal solution exists on $[0,T)$; either $T=\infty$,
or curvature becomes unbounded as $t\nearrow T$ \cite{Ha82}.

\section{Singularities}

Hamilton's compactness, Harnack, and pinching estimates reduce
three-dimensional singularity models to complete nonnegatively curved
limits.  Type-I models have spherical or necklike structure, and a
nonnegatively curved Type-II model is either a positively curved steady
soliton or the product of the cigar soliton with $\mathbb R$
\cite{Ha93, Ha93E, Ha95, Ha95F, Iv}.  Perelman's no-local-collapsing
theorems and local injectivity-radius estimate supply the missing
injectivity control, including Hamilton's Little Loop Lemma
\cite{Ha95F, CCCY}, and yield a uniform canonical-neighborhood theorem
\cite{P1}.  The relevant no-local-collapsing statements are developed
in Theorems 3.3.2 and 3.3.3, and the injectivity-radius consequences in
Theorems 4.2.2 and 4.2.4.

\section{Surgery}

Canonical neighborhoods permit surgery by cutting sufficiently round
necks, attaching caps, and discarding spherical components \cite{Ha97}.
The surgery
parameters are chosen increasingly accurately so that surgery times are
discrete \cite{P2}.  Rescaling arguments use surgical regions pushed away at the
cutoff scale and local smooth extension after surgery; the latter uses
the noncompact uniqueness theorem \cite{P2, CZ05F, CZ05U}.  The
surgery construction is completed in Chapter~8, in particular
Theorem~7.4.3.  Combined with finite extinction for simply connected
manifolds \cite{P3, CM}, it gives the Poincar\'e conjecture.

\section{Long-Time Behavior}

For nonsingular three-dimensional Ricci flows, the long-time geometry
either collapses, converges to constant curvature, or admits a
thick-thin decomposition whose thick part is hyperbolic and whose thin
part collapses; the hyperbolic boundaries are incompressible tori
\cite{Ha99, ScY79, MY, ChG86, ChG90}.  Perelman's analysis extends the
thick-thin approach to flows with surgery.  The thin collapsing result
of Shioya--Yamaguchi for closed manifolds \cite{ShY} suffices for the
geometrization conclusion proved in Theorem~7.7.1.  The explicit
thick-thin statement used here is Theorem~7.6.3; it requires a relation
between its accuracy parameter $\varepsilon$ and collapsing parameter
$w$.

The subsequent chapters develop the Ricci-flow estimates, singularity
models, surgery, and long-time analysis used in these conclusions.
